% Generated draft: requires mathematical and source review.
\documentclass[11pt]{article}
\usepackage[margin=25mm]{geometry}
\usepackage{fontspec}
\setmainfont{Latin Modern Roman}
\usepackage{amsmath,amssymb,mathtools,unicode-math}
\setmathfont{Latin Modern Math}
\usepackage{enumitem,needspace,xurl,hyperref}
\setlength{\parindent}{0pt}
\setlength{\parskip}{5pt}
\setlength{\emergencystretch}{3em}
\newcommand{\N}{\mathbb{N}}
\newcommand{\Z}{\mathbb{Z}}
\newcommand{\Q}{\mathbb{Q}}
\newcommand{\R}{\mathbb{R}}
\newcommand{\C}{\mathbb{C}}
\newcommand{\F}{\mathbb{F}}
\newcommand{\A}{\mathbb{A}}
\newcommand{\PP}{\mathbb{P}}
\newcommand{\E}{\mathbb{E}}
\newcommand{\eps}{\varepsilon}
\newcommand{\dd}{\,\mathrm{d}}
\begin{document}
\section*{UnsolvedMath ID: 30001734}
\textbf{Source problem code:} catalogue entry
\subsection{Definitions and mathematical statement}
The supplied record concerns the local Gross--Prasad conjecture for orthogonal groups. It uses parameter spaces \(\Phi^{\mathrm{orth}}(d)\) and \(\Phi^{\mathrm{orth}}(d-1,\delta)\), a genericity condition, corresponding \(L\)-packets, and a multiplicity \(m(\pi,\sigma)\). The record does not define these objects sufficiently to specify a complete local representation-theoretic setting.

For the purpose of recording the extracted assertion, let \(\mathcal P(\phi,\phi')\) denote the collection of representation pairs described by the source as belonging to the corresponding orthogonal-group \(L\)-packets. This is a placeholder for the source's intended collection; in particular, its distribution over groups or inner forms is unresolved. Write \(m(\pi,\sigma)\) for the source's multiplicity, whose defining representation-theoretic construction is not supplied.

\textbf{Extracted conjectural statement.} For every generic parameter pair
\[
\phi\in\Phi^{\mathrm{orth}}(d),\qquad
\phi'\in\Phi^{\mathrm{orth}}(d-1,\delta),
\]
subject to the admissibility and compatibility conditions intended by the source, there exists a unique pair \((\pi_0,\sigma_0)\in\mathcal P(\phi,\phi')\) such that
\[
m(\pi_0,\sigma_0)=1,
\]
and
\[
m(\pi,\sigma)=0
\quad\text{for every }(\pi,\sigma)\in\mathcal P(\phi,\phi')
\setminus\{(\pi_0,\sigma_0)\}.
\]
Here the admissible values of \(d\) and \(\delta\), the underlying local field, the meanings of the parameter spaces and genericity, the groups and packet conventions, and the definition of multiplicity remain unspecified. These missing data are qualifications on this transcription, not additional hypotheses supplied by the record. The source's genericity wording also requires confirmation as to whether it applies separately to both parameters or to their pair.

The supplied background calls this a multiplicity and epsilon-dichotomy conjecture, but supplies no epsilon-factor formula identifying the distinguished pair. Such a formula cannot be reconstructed from this record alone.
\subsection{Short English statement}
The extracted conjecture says that each generic compatible pair of orthogonal parameters has exactly one representation pair of multiplicity one in the intended packets; every other pair has multiplicity zero. The local setting, packet range, and multiplicity definition need clarification before this becomes a fully specified statement. The supplied label is \texttt{partially\_solved}; this draft does not independently establish that status.
\subsection{Sources}
\begin{itemize}
\item Local UnsolvedMath registry: supplied complete record, permanent ID \(30001734\), exact source problem code \texttt{catalogue entry}, title \emph{Local Gross--Prasad Multiplicity for Orthogonal Groups}, in the Oberwolfach Reports Open Problems collection. The record associates the problem with 2011 and supplies no individual proposer.
\item Supplied original report citation: \emph{Automorphic Forms: New Directions} (2011), Oberwolfach Reports, pp.~719--779, DOI \texttt{10.4171/owr/2011/14}. Supplied literature-source label: original OWR report. \url{https://doi.org/10.4171/owr/2011/14}. The background describes Oberwolfach Reports as documenting mathematical workshops and their open-problem discussions; no exact problem page is supplied.
\item Supplied literature reference, labeled Gan--Gross--Prasad: \url{https://doi.org/10.24033/ast.912}. Further bibliographic details and a precise supporting result are not supplied.
\item Supplied literature reference, labeled local trace formula: \url{https://arxiv.org/abs/2009.13947}. Further bibliographic details and a precise supporting result are not supplied.
\item Supplied literature assessment, checked on \textbf{2026-08-21}: the record labels the problem partially solved and states that major cases, including tempered special orthogonal cases, are proved, while the extracted general formulation is not uniformly settled. This is supplied dated curation derived from the OWR record, not an independently verified literature assessment. The record explicitly requests independent verification of time-sensitive status and cited claims.
\item Supplied statement-verification metadata: \texttt{corrected\_verified}. The accompanying assertion says that the mathematical content was checked against the official Oberwolfach report and that notation, context, and question form were normalized without changing the source-level claim. This is a supplied verification claim, not verification performed for this draft.
\end{itemize}
\end{document}
